\SourceLeaf{400}%
\[
\begin{aligned}
\tfrac{\alpha y'+\beta(y''+y)-\eta}{\alpha x'+\beta(x''+x)-\xi}.\tfrac{y''-y}{x''-x}&=-1\\
\tfrac{\alpha y''+\beta(y+y')-\eta}{\alpha x''+\beta(x+x')-\xi}.\tfrac{y-y'}{x-x'}&=-1
\end{aligned}
\]
Elimination readily gives
\[
\begin{aligned}
\xi&=\tfrac{(y-y')(y'-y'')(y''-y)(\alpha-\beta)}{y(x''-x')+y'(x-x'')+y''(x'-x)}\\
&\quad+\alpha.\tfrac{xy(x''-x')+x'y'(x-x'')+x''y''(x'-x)}{y(x''-x')+y'(x-x'')+y''(x'-x)}\\
&\quad+\beta.\tfrac{y(x''x''-x'x')+y'(xx-x''x'')+y''(x'x'-xx)}{y(x''-x')+y'(x-x'')+y''(x'-x)}
\end{aligned}
\]
\par
The value of \(\eta\) follows from that of \(\xi\) by exchanging all \(x,x',x''\) with the corresponding \(y,y',y''\), as can also be readily seen in advance.

The coordinates of the intersection of the first and last perpendiculars are readily obtained from \(\xi\) and \(\eta\) by interchanging \(x\) with \(x'\) and \(y\) with \(y'\). But this does not change the values of \(\xi\) and \(\eta\), since the coordinates of the points \(A,A',A''\) clearly enter both expressions in the same way. This second intersection therefore clearly coincides with the first, and for the same reason the third intersection also coincides with the first two.

For the centre of gravity, moreover, we evidently have
\[
n=\tfrac23
\]
so
\[
\alpha=\beta=\tfrac13
\]
and therefore
\[
\begin{aligned}\xi&=\tfrac13(x+x'+x'')\\\eta&=\tfrac13(y+y'+y'')\end{aligned}
\]
\par
For the centre of the circumscribed circle we have
\[
n=1
\]
so
\[
\alpha=0,\quad\beta=\tfrac12
\]
\par
For the point of intersection of the perpendiculars from \(A\) and the other vertices themselves, we have
\[
n=0
\]
